\section{Conclusions and Outlook}
\label{sec:conclusions}

Version 4.2 adopts the factor-12 upper length-scale setting for the full 2016
sample and completes the corresponding three-campaign result. The search
intervals and wider GP supports remain 19--90 and 14--135~MeV for 2015, 39--180
and 30--210~MeV for 2016, and 50--250 and 40--300~MeV for 2021. The full 2015
and 2016 samples and the 2021 10\% sample enter a shared-$\eps^2$ likelihood.

The adopted setting follows the controlled factor-8, 10, 12, 15, and 20 range
study. The v4 ceiling is occupied at all 142 full-2016 hypotheses; factor 10
remains occupied at 56, while factors 12, 15, and 20 have no occupied
hypotheses. Factor 12 is the first nonbinding value and is followed by a stable
plateau in log marginal likelihood, observed yield limits, and local $Z$.
Repeated actual fits resolve the reviewed branches, and no GP state, limit, or
$p$-value is interpolated. The rule does not select whichever observed limit is
tighter or whichever local $p_0$ is smaller.

The exact 415-state reconstruction is finite on all 232 combined hypotheses.
For the combined search, 300 shared mass-local conditional background-only
pseudoexperiments provide the median and central limit quantiles. They retain the
reviewed GP states without refitting and use the same asymptotic inner 90\%
\CLs{} construction as the observed limit. The observed-to-median ratio ranges
from 0.514396 at 73~MeV to 3.29400 at 65~MeV. The corresponding mass-local tail
counts reach $0/300$ on the strong side at 71--73~MeV and on the weak side at
21, 65, and 66~MeV; a zero count is finite ensemble resolution, not an exact
probability.

The individual full-2015, full-2016, and 2021 10\% observed limits are reported
alongside the combined result. Auxiliary 100-pseudoexperiment conditional bands
document each standalone likelihood and each pairwise shared-$\eps^2$
combination. They select the first 100 entries from the accepted
300-pseudoexperiment parent stream, reuse campaign draws across scopes, and are
not used as direct-coverage or scan-wide discovery ensembles. The smallest
campaign-level asymptotic $p_0$ values are
$8.46283\times10^{-4}$ at 51~MeV, $3.08781\times10^{-4}$ at 90~MeV, and
$1.05702\times10^{-5}$ at 65~MeV, respectively. The combined minimum is
$3.25918\times10^{-5}$ at 65~MeV ($Z=3.99321$). With
$N_{\mathrm{eff}}=35.3814$, the fixed-card analytic \v{S}id\'ak reference is
$p=0.00115250$ ($Z=3.04783$).

The exact 65~MeV extraction gives
$\widehat{\eps^2}=(6.42706\pm1.61009)\times10^{-6}$. Its dataset-block
decomposition shows that the localized feature is predominantly carried by the
2021 10\% sample, is modestly supported by 2015, and is opposed by 2016 at the
common-coupling yield. The displayed background-subtracted spectra and profiled
residuals document this structure without treating individual bins as independent
significances.

The conditional limit ensemble does not establish direct \CLs{} coverage, refit
the GP, or supply a scan-wise discovery calibration. The analytic \v{S}id\'ak
reference likewise does not account for the observed-data upper-bound study.
Factor 12 is accepted for the v4.2 result, while the post-selection and coverage
qualifications remain explicit.

A ten-toy 2021 source-and-exposure pilot now exercises the paired Poisson generator,
production-geometry GP refit, and hyperparameter export for five declared scenarios:
native 1\%, $1\%\times10$, $1\%\times100$, native 10\%, and $10\%\times10$.
The native-10\% support total is 11.296 times the native-1\% total, so the
$1\%\times10$ versus native-10\% comparison probes source/selection shape and native
normalization rather than a pure exposure scaling. The observed comparisons also move
downward with data fraction: the median $\ell_{\mathrm{opt}}/\sigma_x$ changes from
10.782 to 9.716 between 2016 10\% and full data on the factor-12 card, and from
13.061 to 12.060 between 2021 1\% and 10\% on the factor-15 card. Higher statistics
therefore do not imply a longer optimized scale.

After actual-fit optimizer review, factor 15 truncates the two projected-100\% lanes,
whereas factors 20 and 25 are nonbinding and form the likelihood plateau. Factor 20
is the provisional upper ceiling for the 2021 100\%-statistics projection. If one
ceiling must instead cover all five pilot exposures, the native-1\% boundary contact
at factor 20 makes factor 25 the first universally nonbinding tested value. Both
statements are background-only pilot conclusions, not production-card freezes.

The factor-15-anchored fixed-amplitude screen shows no factor-20 sensitivity loss in
the projected lanes: three fitted-$\sigma_A$ ratios are unity to the quoted precision,
and the alternate native-10\%$\times10$ lane improves by 1.43\%. The four factor-20
paired-response medians lie between 0.971921 and 0.979603 and are essentially unchanged
at factors 15 and 25. The common 2--3\% under-response is therefore not caused by
lifting the upper bound. With ten toys it remains an unresolved closure diagnostic,
not a calibrated bias claim. The pilot itself constructs no expected-limit band or
coverage result, and its unused $q_\mu$ and immutable-path $\eps^2$ outputs are
non-promotable.

The production follow-up should use a predeclared larger ensemble, retain paired toy
identifiers within each truth-source scaling family and across upper-bound settings,
and include independent smooth truth families and injected-signal closure. Its
principal outputs should remain $\ell_{\mathrm{opt}}/\sigma_x$,
$\ell_{\mathrm{opt}}/\ell_{\max}$, boundary occupancy, log marginal likelihood,
repeat stability, background bias, pull, and signal recovery versus exposure. Only
that larger closure study can support a frozen numerical range; direct limit coverage
remains a separate post-freeze requirement.

Production-faithful functional-form closure and direct limit coverage should be
rerun under a card frozen independently of observed limit direction or local
significance. A separate scan-wise
background-only maximum-$q_0$ ensemble would be required for a toy-calibrated global
discovery statement. The eventual full-2021 observed spectrum remains a new analysis
input, not a deterministic scale factor applied to the present 10\% observation.
